\chapter{Resultados}
\label{cap:resultados}

\section{Cobertura e correção funcional}

O conjunto de validação arquivado registra a verificação local bem-sucedida das provas canônicas das oito variantes, com o alvo EVM. A suíte Noir registra 71 funções de teste aprovadas: oito para cada variante inteira, onze para cada variante Field com restrições, onze para a variante com compromisso e três para Field diagnóstico. Esses resultados correspondem aos casos de teste descritos na Seção~\ref{cap:implementacao}.

Na validação externa de geração de testemunhas, os 46 casos inválidos exercitados foram rejeitados: três para cada variante inteira, dez para cada variante Field com restrições e sete para a variante com compromisso. Essas rejeições ocorrem na preparação da testemunha; não representam 46 provas inválidas submetidas a um verificador. O conjunto não inclui casos negativos externos para Field diagnóstico, cuja relação exclui os extremos mas não exige pertencimento ao intervalo inteiro.

A suíte Foundry registra quinze funções de teste aprovadas, abrangendo aceitação das provas válidas e rejeição de alterações de provas ou entradas públicas, entre outras verificações descritas na implementação. Algumas funções percorrem múltiplos casos; por isso, esse total não é somado aos 46 casos de testemunha como se todos representassem a mesma unidade de validação. Os resultados Foundry também não substituem os recibos de transação necessários à comparação final de gas.

Em conjunto, os registros sustentam a aceitação das entradas canônicas e a rejeição dos casos inválidos exercitados nas etapas correspondentes. O alcance dessa evidência é funcional e limitado à suíte executada: não demonstra ausência de toda sub-restrição possível nem segurança formal do backend. A análise de desempenho utiliza essas relações validadas, preservando a distinção entre Field diagnóstico e as variantes de intervalo.

\section{Custos estruturais e verificação EVM}

A Tabela~\ref{tab:circuitos} apresenta as contagens estruturais e os tamanhos das provas canônicas validadas. As três variantes inteiras possuem a mesma contagem ACIR, enquanto a contagem de gates difere. Nas variantes Field com restrições, a contagem ACIR também permanece constante entre larguras. Assim, a igualdade nessa contagem intermediária não implica igualdade de estrutura no backend.

\begin{table}[htbp]
\centering
\caption{Operações e tamanho das provas das oito variantes}
\label{tab:circuitos}
\begin{tabular}{lrrr}
\hline
Variante & ACIR & Gates & Prova (bytes) \\
\hline
Field diagnóstico & 5 & 40 & 4\,928 \\
u8 & 11 & 93 & 5\,312 \\
u32 & 11 & 2\,791 & 7\,232 \\
u64 & 11 & 2\,838 & 7\,232 \\
Field restrito 8 & 7 & 91 & 5\,312 \\
Field restrito 32 & 7 & 2\,789 & 7\,232 \\
Field restrito 64 & 7 & 2\,836 & 7\,232 \\
u64 com compromisso & 37 & 14\,348 & 8\,000 \\
\hline
\end{tabular}
\par\smallskip
\begin{minipage}{\textwidth}
\renewcommand{\baselinestretch}{1}\small
Fonte: elaboração própria a partir dos artefatos validados do experimento (2026).\par
Nota: ACIR e gates são contagens de camadas distintas. Provas canônicas com alvo EVM; entradas e versões descritas na Seção~\ref{cap:implementacao}. Field diagnóstico não verifica pertencimento ao intervalo.
\end{minipage}
\end{table}


As provas das variantes de 32 e 64 bits têm o mesmo tamanho no conjunto apresentado, tanto para inteiros quanto para Field com restrições. A variante com compromisso possui mais operações ACIR, mais gates e uma prova maior que a variante u64 sem compromisso. Essas observações descrevem os artefatos desta configuração; a comparação de latência depende da coleta correspondente.

A Tabela~\ref{tab:verificadores} apresenta os tamanhos dos modelos compilados de runtime dos verificadores. Todos ficam abaixo de 24.576 bytes, com margens de 330 a 333 bytes. A margem é uma diferença aritmética em relação ao limite EIP-170. Na rodada de gas, os oito verificadores foram implantados e seus códigos conferidos, incluindo a vinculação das bibliotecas e os valores inseridos pelo construtor.

\begin{table}[htbp]
\centering
\caption{Tamanho dos modelos compilados de código dos verificadores}
\label{tab:verificadores}
\begin{tabular}{lrr}
\hline
Variante & Runtime (bytes) & Margem (bytes) \\
\hline
Field diagnóstico & 24\,243 & 333 \\
u8 & 24\,246 & 330 \\
u32 & 24\,244 & 332 \\
u64 & 24\,245 & 331 \\
Field restrito 8 & 24\,245 & 331 \\
Field restrito 32 & 24\,243 & 333 \\
Field restrito 64 & 24\,245 & 331 \\
u64 com compromisso & 24\,245 & 331 \\
\hline
\end{tabular}
\par\smallskip
\begin{minipage}{\textwidth}
\renewcommand{\baselinestretch}{1}\small
Fonte: elaboração própria a partir dos artefatos validados do experimento (2026).\par
Nota: Margem calculada em relação a 24.576 bytes (EIP-170). Os tamanhos correspondem aos modelos compilados, incluindo posições de vinculação de bibliotecas; não substituem a conferência do código efetivamente implantado. Não são medidas de gas.
\end{minipage}
\end{table}


A Tabela~\ref{tab:gas} apresenta o gas observado nas 40 provas selecionadas. As 40 chamadas de verificação retornaram verdadeiro e os 80 recibos registraram sucesso. Em cada prova, as duas transações tiveram gas total idêntico. O piso de calldata não vinculou em nenhuma dessas verificações, permitindo recuperar a execução líquida em todos os casos.

\begin{table}[htbp]
\centering
\caption{Gas de verificação das oito variantes no fork local de Sepolia}
\label{tab:gas}
\small
\begin{tabular}{lrrrr}
\hline
 & \multicolumn{3}{c}{Gas total} & Execução \\
Variante & Mínimo & Mediana & Máximo & líquida \\
\hline
Field diagnóstico & 1\,935\,842 & 1\,935\,926 & 1\,935\,926 & 1\,838\,390 \\
u8 & 2\,003\,208 & 2\,003\,292 & 2\,003\,400 & 1\,899\,684 \\
u32 & 2\,340\,880 & 2\,340\,904 & 2\,341\,012 & 2\,206\,708 \\
u64 & 2\,340\,796 & 2\,340\,892 & 2\,340\,964 & 2\,206\,708 \\
Field restrito 8 & 2\,003\,196 & 2\,003\,292 & 2\,003\,352 & 1\,899\,684 \\
Field restrito 32 & 2\,340\,832 & 2\,340\,892 & 2\,340\,952 & 2\,206\,708 \\
Field restrito 64 & 2\,340\,904 & 2\,340\,916 & 2\,340\,964 & 2\,206\,708 \\
u64 com compromisso & 2\,479\,655 & 2\,479\,739 & 2\,479\,811 & 2\,332\,639 \\
\hline
\end{tabular}
\par\smallskip
\begin{minipage}{\textwidth}
\renewcommand{\baselinestretch}{1}\small
Fonte: elaboração própria a partir dos recibos da coleta de 26/10/2026.\par
Nota: valores em unidades de gas. Resumo de cinco provas por variante; duas transações por prova, com custos idênticos entre repetições. A execução líquida foi constante nas cinco provas. Fork no bloco 11.999.881, Anvil 1.1.0 com Prague; solc 0.8.28, alvo Cancun e parâmetro optimizer\_runs=1. Field diagnóstico possui relação distinta das variantes de intervalo.
\end{minipage}
\end{table}




Nas cinco provas de cada variante, a execução líquida permaneceu constante, enquanto o gas total variou com a composição da calldata. Os pares Field restrito 8/u8, Field restrito 32/u32 e Field restrito 64/u64 apresentaram os mesmos valores de execução líquida. As variantes de 32 e 64 bits também coincidiram nessa métrica. Essas igualdades descrevem as provas selecionadas e a implementação fixada; não decorrem de um teste de equivalência sobre uma população de provas.

A variante u64 com compromisso apresentou maior custo de execução que u64 sem compromisso. Essa diferença abrange a verificação da relação adicional e de suas entradas públicas, não o custo isolado de uma operação de hash. Field diagnóstico apresentou o menor valor observado, mas verifica um predicado distinto e, portanto, não representa uma alternativa semanticamente equivalente de prova de intervalo.

Os custos de implantação estão separados na Tabela~\ref{tab:implantacao}. Todos os 16 recibos registraram sucesso. A soma do verificador e de sua biblioteca variou de 6.670.986 a 6.671.659 gas entre as variantes. A rodada implantou uma biblioteca para cada verificador; os valores não pressupõem compartilhamento ou reaproveitamento de uma implantação anterior.

\begin{table}[htbp]
\centering
\caption{Gas total de implantação dos verificadores e bibliotecas vinculadas}
\label{tab:implantacao}
\begin{tabular}{lrrr}
\hline
Variante & Verificador & Biblioteca & Soma \\
\hline
Field diagnóstico & 5\,294\,220 & 1\,376\,766 & 6\,670\,986 \\
u8 & 5\,294\,893 & 1\,376\,766 & 6\,671\,659 \\
u32 & 5\,294\,536 & 1\,376\,766 & 6\,671\,302 \\
u64 & 5\,294\,632 & 1\,376\,766 & 6\,671\,398 \\
Field restrito 8 & 5\,294\,629 & 1\,376\,766 & 6\,671\,395 \\
Field restrito 32 & 5\,294\,260 & 1\,376\,754 & 6\,671\,014 \\
Field restrito 64 & 5\,294\,608 & 1\,376\,766 & 6\,671\,374 \\
u64 com compromisso & 5\,294\,608 & 1\,376\,766 & 6\,671\,374 \\
\hline
\end{tabular}
\par\smallskip
\begin{minipage}{\textwidth}
\renewcommand{\baselinestretch}{1}\small
Fonte: elaboração própria a partir dos recibos da rodada de 26/10/2026.\par
Nota: valores em unidades de gas; uma transação por verificador (CREATE) e por biblioteca (CREATE2). Soma dos dois recibos associados à variante, sem amortização ou custo de infraestrutura previamente disponível no fork. Mesma configuração da Tabela~\ref{tab:gas}.
\end{minipage}
\end{table}


\section{Execução da coleta e incidentes}
Até 11 de novembro de 2026, foram concluídos dez pares diários: de 26 de outubro a 1º de novembro e nos dias 3, 6 e 11 de novembro. Os vinte registros de sessão desses pares totalizam 18.400 provas verificadas, sendo 2.400 de aquecimento e 16.000 medidas, sem falhas de geração ou verificação registradas nessas sessões completas. Esse total descreve os pares retidos para a análise; tentativas interrompidas permanecem arquivadas separadamente. A análise final de latência depende dos dez pares restantes e da avaliação de elegibilidade e dos pressupostos.

Em 29 de outubro, uma primeira tentativa completou apenas a sessão nativa e foi interrompida antes do WASM. Uma repetição integral foi autorizada antes de sua execução; o par completo dessa segunda tentativa foi retido, preservando o registro anterior sem combinar suas sessões. Em 3 de novembro, uma interpretação de reinício levou à reutilização das sementes do primeiro par. O autor esclareceu posteriormente a intenção de continuidade, e o par foi preservado com sua proveniência original. Essas emendas ocorreram durante a coleta e devem acompanhar a avaliação final, sem serem apresentadas como parte inalterada da alocação inicial.

Não há pares coletados em 2, 4, 5 e de 7 a 10 de novembro. Em 6 de novembro, um lançamento foi encerrado durante a pausa inicial, antes de qualquer tentativa, e preservado; a coleta posteriormente autorizada nessa data foi concluída. Em 11 de novembro, o controle de espaço em disco bloqueou a execução antes das medições; após liberação de espaço e novo preflight, o par foi executado integralmente. Ausências e bloqueios anteriores a tentativas não são contados como falhas de prova.

Na preparação da rodada de gas, duas tentativas foram interrompidas em consultas de rede. Outras duas chegaram às implantações, mas pararam na conferência do código: primeiro pela omissão de bibliotecas implantadas via CREATE2 no inventário e depois pelo tratamento incompleto dos valores imutáveis inseridos pelo construtor. Os controles foram corrigidos e os registros anteriores preservados. As 80 verificações analisadas pertencem à rodada concluída após essas correções, mantendo as provas selecionadas e o mesmo bloco e hash fixados antes das implantações. O diagnóstico separado do piso resultou em 61.000 gas, conforme esperado para o vetor descrito no método.

\section{Latência e contrastes diários}
\pendente{Inserir as estimativas dos 16 contrastes e seus intervalos após completar e validar a alocação de vinte pares. Esta versão não apresenta inferência parcial de latência.}
